\clearpage
\fancyhead[L]{\small ANWENDUNG DER WAHRSCHEINLICHKEITSRECHNUNG}
\pdfbookmark[0]{Nachlass. Anwendung der Wahrscheinlichkeitsrechnung auf die Bestimmung der Bilanz für Witwenkassen}{witwen}
\SourceLeaf{124}
\begin{center}
{\Large NACHLASS.\par}\bigskip
{\large [ANWENDUNG DER WAHRSCHEINLICHKEITSRECHNUNG\par}
\medskip{\Large AUF DIE BESTIMMUNG DER BILANZ\par}\medskip
{\large FÜR WITWENKASSEN ]\par}
\medskip\rule{.15\linewidth}{.3pt}\par\bigskip
{\large [I.]\par}
\pdfbookmark[1]{I. Allgemeine Übersicht über die Methode}{witwen-I}
\textit{[Allgemeine Übersicht über die Methode.]}\par
\textit{[Auszug aus einem Votum bei der schriftlichen Abstimmung im Universitätssenat.]}\par
\end{center}
Das vorstehende [hier eingeklammerte] Votum von Herrn Prof. D.: [Das Königliche Universitätskuratorium scheint zu befürchten, dass die Kasse bei der großen Zahl der jetzt vorhandenen Witwen über kurz oder lang nicht mehr imstande sein werde, die inzwischen auf 250 Thl. angewachsenen Pensionen zu bestreiten. Es verlangt daher einen Bericht darüber, ob es einen begründeten Anlass zu einer solchen Besorgnis gibt und mit welchen Mitteln die möglicherweise drohende Gefahr abgewendet werden könne. .\,.\,.\,] trifft den eigentlichen \textit{Kern der Frage} so genau, dass ich der ersten Hälfte dieses Votums nur wörtlich zustimmen kann. Wenn bezweifelt wird, ob die Kasse der ihr obliegenden Verpflichtung dauerhaft nachkommen kann, dann ist dies wahrlich der ungeeignetste Zeitpunkt, \textit{größere} Anforderungen an sie zu stellen.

\indent Ich kann mich noch an die öffentliche Meinung über diese Anstalt vor bis zu 40 Jahren erinnern. Schon damals galt sie als ein herrliches Kleinod der Universität, einzigartig in ihrer Art, und zwar gerade wegen ihrer Besonderheiten. \textit{Völlige} Freiheit, ob man beitreten wollte oder nicht, ja sogar mit einem vergleichsweise geringen Opfer wieder einzutreten, nachdem man ausgetreten war; ein sehr geringer Beitrag. Und damals betrug die Pension nur 150 oder 160 Thl. Nicht die Höhe der Pension war das Anziehende, sondern die \textit{großzügige Art, wie} demjenigen, der Göttinger Professor werden konnte, eine sichere Unterstützung für eine hinterbliebene Witwe angeboten wurde, mit der Aussicht, sie unter weiser Verwaltung nach und nach noch erhöht zu erhalten. Wer mehr wünschte, beteiligte sich daneben noch an einer anderen Witwenkasse. Nun ist die Pension auf 250 Thl. gestiegen, und die großzügige Art ist bis heute dieselbe geblieben. Gebe Gott, dass es niemals nötig werde, daran \textit{irgendetwas} zu ändern! Zwangsweise prozentuale Abzüge vom Gehalt, um durch Drehen am Stundenzeiger das zu erreichen, was nur das allmähliche Vorrücken des Minutenzeigers gewähren kann, wären nicht nur viel zu unwirksam, um den Zweck zu erreichen, sondern würden den edlen Charakter der Anstalt völlig zerstören.

\par\SourceLeaf{125}
Ich bin daher der Meinung, sämtliche Änderungsvorschläge von Herrn Universitätsrat O. für den Augenblick ganz auf sich beruhen zu lassen; in dem P. M., für das wir ihm zu lebhaftem Dank verpflichtet sind, ist gezeigt worden, dass keine unmittelbare Gefahr besteht. Selbst wenn in den nächsten Jahren durch weitere hinzukommende Witwen an die Stelle des letzten, noch immer erfreulichen Überschusses ein gewisses Defizit treten sollte, darf man nicht vergessen, dass die angesammelten Überschüsse zum Teil gerade dazu bestimmt sind, solche durch vorübergehende Umstände entstandenen Schwankungen aufzufangen.

Eine \textit{gründliche} Untersuchung halte ich aber in Übereinstimmung mit dem Reskript und den von Seiner Magnifizenz geäußerten Ansichten durchaus für notwendig. Selbst bei der zuversichtlichsten Einschätzung des Zustands der Gesellschaft wird eine solche Untersuchung jedenfalls zumindest später notwendig werden müssen, schon aus folgendem Grund.

Wenn ich meine \textit{Meinung} äußern darf, \textit{bevor} eine gründliche, auf strenger Berechnung beruhende Untersuchung stattgefunden hat, glaube ich, dass die derzeit große Zahl der Witwen als anomal betrachtet werden muss. Es ist wahr, dass die Zahl der teilnehmenden Professoren in einem gewissen Verhältnis zur Zahl der Witwen stehen muss und dass die erstere Zahl jetzt viel größer ist als früher. Die \textit{jetzt} hohe Zahl der Witwen steht damit aber in keinerlei Zusammenhang. Bleibt die Zahl der teilnehmenden Professoren von nun an stets so groß, dann ist dies \textit{ein sehr ernsthaft zu erwägender Punkt}, aber wegen der fernen Zukunft, nicht wegen der Gegenwart; erst nach 20 oder 30 und mehr Jahren können die Folgen \textit{davon} sehr deutlich werden.

Unter dieser Voraussetzung ist mit Wahrscheinlichkeit anzunehmen, dass die Zahl der Witwen vielleicht schon nach wenigen Jahren wieder abnehmen, vielleicht erheblich abnehmen wird und dass daher die angesammelten Überschüsse von 8149 Thl. wieder anwachsen, vielleicht erheblich anwachsen werden. \textit{Ob aber die Zahl der Witwen beispielsweise innerhalb von \textup{10} Jahren auf oder unter \textup{15} abnehmen wird, ist viel ungewisser.} Angenommen, die Überschüsse wären auf 15000 Thl. oder mehr angewachsen, die Zahl der Witwen bliebe aber hartnäckig bei 16 stehen: Was soll dann geschehen? Einerseits will man die Überschüsse nicht unbegrenzt anwachsen lassen, andererseits steht das \textit{Statut} einer Erhöhung der Pension entgegen. Dann \textit{muss} gründlich geprüft werden, ob und \textit{in welchem Maß} man das \textit{Statut ändern} darf, ohne die Gesellschaft zu gefährden.

Dem Vertrauen, mit dem Seine Magnifizenz und einige der Herren Kollegen mich ehren, indem sie wünschen, dass ich eine solche Prüfung übernehme, will ich mich nicht entziehen. Dabei soll nämlich eine auf Sterblichkeitsgesetzen und der Wahrscheinlichkeitsrechnung beruhende Bilanz zwischen dem Vermögen der Anstalt und ihren Verpflichtungen gezogen werden. Ich muss jedoch Folgendes vorausschicken.

Erstens haben diejenigen Herren eine sehr falsche Vorstellung von der Langwierigkeit solcher Rechnungen, die glauben, sie könnten innerhalb von vier Wochen abgeschlossen werden. Ich kann mich daher umso weniger auf eine bestimmte Frist festlegen, je kleiner der Teil meiner Zeit sein wird, den ich darauf verwenden kann.

Zweitens lassen sich die Rechnungen ohne die nötigen \textit{Daten}, von denen derzeit \textit{überhaupt nichts} vorliegt, gar nicht gründlich durchführen. Worin die erforderlichen Daten bestehen, werde ich weiter unten angeben; ohne sie kann ich mich auf gar nichts einlassen. Ob, auf welche Weise und wie bald sie sich zusammenbringen lassen, muss ich ganz der Kirchendeputation überlassen.

Drittens enthält die besondere Einrichtung unserer Witwenkasse mehrere Elemente, die vom Sterblichkeitsgesetz unabhängig sind und sich einer Berechnung nicht unterwerfen lassen. Deshalb wird dem Endergebnis notwendigerweise eine gewisse Unvollkommenheit anhaften; ich hoffe jedoch, dass sich Ersatzmittel finden lassen, durch deren Verwendung diese Unvollkommenheit unerheblich sein wird.

Ich will nun versuchen, kurz anzudeuten, worauf es bei dieser Arbeit ankommt.

\par\SourceLeaf{126}
Die Verpflichtungen der W. K. zerfallen in drei Hauptrubriken.

I. Verpflichtungen gegenüber den jetzt vorhandenen Witwen und gegebenenfalls deren Kindern,

II. Verpflichtungen gegenüber denjenigen Professoren, die jetzt Teilnehmer der W. K.-Gesellschaft sind.

III. Verpflichtungen gegenüber den künftig beitretenden Professoren.

Ad I. Die Verpflichtung gegenüber jeder einzelnen Witwe ohne minderjährige Kinder hat genau den Wert einer Leibrente für sie und lässt sich daher genau berechnen, wenn man ihr gegenwärtiges Alter kennt und für Sterbetafel und Zinssatz eine bestimmte Wahl trifft. Dass ein Unternehmer, der diese Verpflichtung zu diesem Preis übernähme, in jedem einzelnen Fall eine Art Glücksspiel spielt, versteht sich von selbst (ich werde daher im Folgenden, wo dieser Hinweis immer wieder nötig wäre, darauf verzichten, da dies jedem mit einiger mathematischer Bildung bekannt ist). Ein Unternehmer aber, der mit einer \textit{sehr großen Zahl} von Personen denselben Vertrag abschlösse, hätte bei richtiger Berechnung mit moralischer Gewissheit nur Schwankungen zu erwarten, die im Verhältnis zum Ganzen unerheblich sind.

Sind Kinder vorhanden, muss die Rechnung abgeändert werden, behält aber dieselbe Gültigkeit wie im vorigen Fall. Natürlich muss auch das Alter der Kinder bekannt sein.

Schließlich wäre bei unserer Witwenkasse, die sich wiederverheiratende Witwen ausschließt, streng genommen noch eine weitere Abänderung nötig, die scheinbar\footnote{Es gehört nicht hierher, zu begründen, warum ich diese Einrichtung nur für \textit{scheinbar} vorteilhaft halte. Ich bin aber gern bereit, jedem, der sich dafür interessiert, die Gründe zu nennen.} zum Vorteil der Witwenkasse ist, sich aber natürlich nicht im Voraus berechnen lässt und jedenfalls praktisch \(=0\) zu setzen ist.

Damit lässt sich auch der Gesamtbetrag von I. in Geld bewerten; einen diesem Betrag gleichen Teil des Kapitalvermögens der Kasse muss man als dadurch gebunden betrachten.

Ad II. Ebenso ließe sich auch die Verpflichtung gegenüber jedem einzelnen Professor, der jetzt verheiratetes Mitglied der Gesellschaft ist, in Geld bewerten, wenn es in unserer Gesellschaft genauso wäre wie in denjenigen freiwilligen Gesellschaften, bei denen der jährliche Beitrag oder das Eintrittsgeld nach dem Alter des eintretenden Ehepaars festgelegt wird. Eine solche Gesellschaft unterscheidet sich von unserer aber in Folgendem.

A. Bei jener erlischt der Vertrag, wenn die Frau vor dem Mann stirbt. Soll er bei einer Wiederverheiratung erneuert werden, ist es ein ganz neuer Vertrag, der nach dem Alter der zweiten Frau festzulegen ist. Bei uns gibt es auch unverheiratete Mitglieder, die eine Braut beliebigen Alters wählen können, und ebenso Witwer, die sich möglicherweise wieder verheiraten. All dies lässt sich jetzt aber gar nicht im Voraus berechnen. Ich würde jedoch annehmen, dass es sich ungefähr ausgleichen würde, wenn man für die jetzt verheirateten Mitglieder nach ihrem eigenen Alter und dem ihrer gegenwärtigen Frauen rechnete und dann für die übrigen, derzeit nicht verheirateten Mitglieder den Mittelwert jener ersten Ergebnisse zugrunde legte. Möglicherweise werden bei einigen jetzt verheirateten Mitgliedern einmal nicht ihre gegenwärtigen Frauen, sondern zweite oder dritte Frauen die Witwenpension beziehen; dagegen wird aber zweifellos ein Teil der jetzt nicht verheirateten Mitglieder in diesem Stand bleiben. Ich sehe zumindest nicht, was man hier mehr tun könnte, als die beiden Arten von Möglichkeiten, von denen die einen zum Nachteil und die anderen zum Vorteil der Kasse sind, sich gegenseitig aufheben zu lassen.

B. Außerdem besteht noch der Unterschied, dass Kinder, die vielleicht jetzt noch gar nicht geboren sind, später möglicherweise an den Leistungen teilhaben. Auch das lässt sich daher auf diese Weise nicht veranschlagen; ich glaube jedoch, dass sich für diese Unvollkommenheit ein völlig ausreichender Ersatz finden lässt, den ich hier aber nicht näher entwickeln will, um nicht allzu ausführlich zu werden.

